\documentclass[11pt]{article}
\usepackage[margin=1in]{geometry}
\usepackage[T1]{fontenc}
\usepackage[utf8]{inputenc}
\usepackage{lmodern}
\usepackage{amsmath,amssymb,amsthm,mathtools}
\usepackage{enumitem}
\usepackage{booktabs,longtable,array}
\usepackage{xurl}
\usepackage[hidelinks]{hyperref}
\usepackage{xr-hyper}
\usepackage[backend=biber,style=alphabetic,maxbibnames=8]{biblatex}
\externaldocument[mono-]{../../monolithic/QICN_MONOLITHIC}
\addbibresource{../../release/references.bib}
\setlength{\emergencystretch}{3em}

\newtheorem{definition}{Definition}[section]
\newtheorem{theorem}[definition]{Theorem}
\newtheorem{proposition}[definition]{Proposition}
\newtheorem{lemma}[definition]{Lemma}
\newtheorem{corollary}[definition]{Corollary}
\newtheorem{assumption}[definition]{Assumption}
\newtheorem{remark}[definition]{Remark}
\newtheorem{criterion}[definition]{Criterion}

\newcommand{\Hist}{\operatorname{Hist}}
\newcommand{\Resp}{\operatorname{Resp}}
\newcommand{\QGauge}{\operatorname{Gauge}}
\newcommand{\Iint}{I_{\mathrm{int}}}
\newcommand{\ThetaS}{\Theta_S}
\newcommand{\VTheta}{V_{\Theta}}
\newcommand{\VR}{V_R}
\newcommand{\Cuts}{\operatorname{Cuts}}
\newcommand{\Comp}{\operatorname{Comp}}
\newcommand{\Cert}{\mathsf{Cert}}
\newcommand{\Audit}{\mathsf{Audit}}

\title{QICN v22 Internal Note: Independent Finite Separator-Complete Incidence Package and Conditional Closure of \texorpdfstring{$I_{\mathrm{int}}$}{I_int}}
\author{QICN internal formal-audit scaffold}
\date{2026-05-27}

\begin{document}
\maketitle

\section*{Governance boundary}
This document does not certify external support, consciousness, phenomenality, identity transfer, or bridge-burden closure. It is an internal mathematical and implementation scaffold. It does not report human-subject data, external adjudication, reviewer certification, or empirical validation. The package constructed in this note is finite and synthetic; it may be used as a conditional formal witness only after its catalog-completeness premise is independently accepted.

\section{Scope and layer separation}

\begin{description}[leftmargin=2.9cm]
  \item[Ontology.] A system $S$ is treated only as an admissible object with histories, identity data, intervention-response coordinates, and separator tests. No consciousness, agency, subjecthood, or phenomenal identity is attributed.
  \item[Mathematical model.] The relevant structure is a finite bipartite incidence graph $G_S=(\VTheta\sqcup\VR,E)$ plus an explicit finite factorization-cut universe.
  \item[Implementation.] The finite package is represented by \path{docs/theory/FINITE_SEPARATOR_COMPLETE_PACKAGE_v22.json}; it is audited by \path{scripts/audit-finite-separator-package.js}. The product-separator negative control is represented separately.
  \item[Language.] ``Separator-complete'', ``response-operative'', and ``connected'' are formal predicates over the finite package. They are not synonyms for consciousness or identity transfer.
  \item[Interpretation.] The result is aggressive but bounded: accepted finite separator-complete connected incidence implies atomicity relative to the declared finite factorization universe, and therefore conditional closure of $\Iint$. It is not a global proof from rigid identity plus continuity plus fidelity alone.
\end{description}

\section{Corpus anchors}
This note extends the v20 closure note and is anchored to the monolithic Paper 5 claims. Proposition~\ref{mono-06-paper5:prop:integration-transfer} is the original integration-transfer anchor, and Theorem~\ref{mono-06-paper5:thm:iint-faithful-factorization-triviality} is the prior conditional theorem under an atomic separator. The formal style follows the framework's Paper 5 surface~\cite{paper5}, categorical factorization language~\cite{maclane1998}, and finite graph reasoning compatible with causal graph vocabulary~\cite{pearl2009}. These references orient the formalism only; they do not externally validate QICN.

\section{Finite package primitives}

\begin{definition}[Finite response-separator incidence package]
A finite response-separator incidence package for a system presentation $S$ is a tuple
\[
  \Cert(S)=(\VTheta,\VR,E,\mathcal P,\mathcal E,\mathcal C),
\]
where:
\begin{enumerate}[label=(\roman*)]
  \item $\VTheta$ is a finite set of elementary separator vertices declared by a separator catalog $\mathcal C$;
  \item $\VR$ is a finite set of elementary intervention-response coordinate vertices;
  \item $E\subseteq \VTheta\times\VR$ is the incidence relation;
  \item $\mathcal P$ is a finite set of typed perturbation records witnessing the edges of $E$;
  \item $\mathcal E$ is an enumeration policy for non-trivial binary factorization cuts over $\VTheta\sqcup\VR$.
\end{enumerate}
An edge $(\theta,r)\in E$ is admissible only if its perturbation record preserves the declared identity typing and history typing while producing a positive separator-margin change.
\end{definition}

\begin{definition}[Independent construction boundary]
The package is independently constructed when $\VTheta$, $\VR$, $E$, and $\mathcal P$ are declared before applying the atomicity predicate, and the construction rule does not query whether $\Theta_S$ is atomic. Operationally, the v22 package is generated from a literal separator catalog, response-coordinate catalog, and perturbation table. The audit then checks graph connectedness and cut obstruction afterward.
\end{definition}

\begin{definition}[Finite separator-completeness]
A package $\Cert(S)$ is finite separator-complete relative to its declared universe when:
\begin{enumerate}[label=(\roman*)]
  \item every declared separator in the catalog $\mathcal C$ appears exactly as a vertex of $\VTheta$;
  \item every $\theta\in\VTheta$ is response-operative, meaning it has at least one incident edge witnessed by a positive typed perturbation record;
  \item every $r\in\VR$ has at least one incident edge;
  \item every admissible factorization considered by the finite model induces a binary cut listed by the enumeration policy $\mathcal E$;
  \item a factor-local decomposition corresponds to a non-trivial cut with no incidence edge crossing the cut.
\end{enumerate}
This is a finite relative condition. It does not claim completeness over undeclared separators, unrecorded response coordinates, infinite systems, or external empirical systems.
\end{definition}

\begin{definition}[Connected finite incidence]
The incidence graph of $\Cert(S)$ is connected when the undirected bipartite graph $G_S=(\VTheta\sqcup\VR,E)$ has exactly one connected component after rejecting isolated non-operative vertices.
\end{definition}

\begin{definition}[Relative atomicity]
The separator family is atomic relative to $\Cert(S)$ when no non-trivial finite factorization cut enumerated by $\mathcal E$ is factor-local, i.e. every enumerated non-trivial cut has at least one crossing incidence edge.
\end{definition}

\begin{definition}[Hash-stable package integrity]
A finite package is hash-stable when the declared package digest equals
\[
  \operatorname{SHA256}(\operatorname{stable\_json}(P\setminus\{\texttt{package\_sha256}\})).
\]
The digest excludes only the top-level digest field itself. Declared audit metadata, admission rules, perturbation records, and limitations remain inside the digest. This avoids the self-referential hash defect in which recomputing a digest over a package that already contains its own digest changes the object being hashed.
\end{definition}


\section{The constructed v22 package}

\begin{criterion}[v22 connected finite package]
The file \path{FINITE_SEPARATOR_COMPLETE_PACKAGE_v22.json} declares six separator vertices, five response-coordinate vertices, twelve typed incidence edges, and twelve typed perturbation records. The audit enumerates all canonical binary cuts over the eleven vertices, restricts to non-trivial factor cuts containing at least one separator and one response on each side, and verifies that no such cut is edge-closed. It also stores the recomputed counts inside the package as declared audit metadata and verifies that the package hash is stable under the top-level digest-exclusion rule.
\end{criterion}

\begin{criterion}[Product-separator negative control]
The file \path{FINITE_SEPARATOR_COMPLETE_PACKAGE_PRODUCT_NEGATIVE_CONTROL_v22.json} deliberately declares two disconnected product components. The same audit must reject it as an atomicity witness. This prevents the validator from merely accepting any syntactically complete finite package.
\end{criterion}

\section{Main theorem}

\begin{theorem}[Finite separator-complete connected incidence implies relative atomicity]
Let $\Cert(S)$ be a finite separator-complete response-separator incidence package. If its incidence graph is connected, then $\Theta_S$ is atomic relative to the finite factorization universe declared by $\Cert(S)$.
\end{theorem}

\begin{proof}
Assume for contradiction that $\Theta_S$ is not atomic relative to $\Cert(S)$. Then there exists an enumerated non-trivial factorization cut
\[
  \VTheta\sqcup\VR = A\sqcup B
\]
with both $A$ and $B$ containing at least one response-operative separator and at least one response-coordinate vertex, and with no edge of $E$ crossing from $A$ to $B$.

Because the cut is non-trivial and all represented separators and response vertices are operative, both induced subgraphs are non-empty. The absence of crossing edges means that the graph decomposes as the disjoint union of the subgraph induced by $A$ and the subgraph induced by $B$. Hence $G_S$ has at least two connected components. This contradicts connectedness. Therefore no enumerated non-trivial factorization cut is factor-local, which is exactly relative atomicity for the finite package.
\end{proof}

\begin{corollary}[v22 finite-witness conditional closure of $\Iint$]
If the v22 finite package is accepted as separator-complete for the declared finite system model, if its audit passes, and if the factorization-loss terms are faithful to history, identity, response, and gauge preservation, then $\Iint(S)$ exists and is unique as a gauge-invariant binary factorization-gap invariant relative to that finite model.
\end{corollary}

\begin{proof}
The audit-passing finite package gives connected finite incidence. The previous theorem gives relative atomicity. The prior conditional theorem under atomic separator, anchored by Theorem~\ref{mono-06-paper5:thm:iint-faithful-factorization-triviality}, then applies to the finite model: a zero-loss non-trivial factorization would require a factor-local separator cut, but relative atomicity forbids such a cut. Therefore the factorization-gap predicate defining $\Iint$ is well-defined and invariant under admissible gauge transformations that preserve the finite package. Uniqueness follows from equality of the binary zero-set for zero-loss factorizations.
\end{proof}

\section{Non-circularity and reviewer burden}

\begin{remark}[Reviewer burden]
A reviewer must verify that the separator catalog and response-coordinate catalog can be obtained operationally without first proving atomicity. The admissible workflow is: declare elementary separator tests, declare elementary response coordinates, record typed perturbation evidence, construct the incidence graph, enumerate cuts, and only then apply connectedness. The inadmissible workflow is: assume no factor-local separator partition exists and then build the catalog around that assumption. If the inadmissible workflow is used, the v22 certificate is circular and must be rejected.
\end{remark}

\begin{remark}[What this closes]
The v22 construction closes a finite, implementation-checkable route: accepted finite separator-complete connected incidence implies relative atomicity, which gives conditional $\Iint$ closure. It does not close the global atomic-separator lemma for arbitrary continuous or infinite systems, and it does not transform synthetic internal evidence into empirical validation.
\end{remark}

\begin{remark}[Negative control interpretation]
The product-separator negative control fails the connectedness and factor-local cut audit. This is the expected behavior. It shows that the validator distinguishes a connected finite witness from a product-decomposed incidence graph. It does not show that real systems satisfy the positive package.
\end{remark}

\section{Implementation status}
The implementation consists of the package constructor, validator, positive package, negative control package, independent catalog protocol, blinding protocol, double-blind review packet, and generated audit report. The validator is deterministic and computes a package hash that excludes only the top-level digest field. The full result is recorded in \path{docs/reports/FINITE_SEPARATOR_COMPLETE_PACKAGE_v22_AUDIT.json}. Promotion to a source paper or theorem registry entry still requires human mathematical curation and an explicit decision about the finite-universe premise.

\printbibliography

\end{document}
